\section{Conclusion}\label{sec:conclusion}

\bfs 0.1.26 passed all \xfsXRan{} xfstests tests that ran on x86-64
and \xfsAPass{} of the \xfsARan{} that ran on aarch64, out of
\xfsTotal{}; xfstests declined the others, which are untested. Under
emulated single-bit upsets on the read path, it never returned wrong
data: it returned the correct bytes after every
flip that reached a data block and refused the one read it could not
vouch for. In the same
conditions ext4 and ext3 returned wrong data without an error and btrfs
refused. These results are bounded by their load, at most
\kernMaxPerDecode{} erroneous symbols per codeword and per decode against
eight correctable; the correction limit, bursts and miscorrection remain
to be measured, with \texttt{DATA\_CSUM} enabled.

Two findings matter beyond the numbers. The injector's record was wrong
in a way the outcomes alone would never have shown; checking it against
the kernel's own account of each correction found the error, and its
cause in the kernel source. And the indirect-block parity, as built,
detects without repairing, which leaves a pointer flipped onto another
allocated block undetected on volumes without \texttt{DATA\_SELFID}.
Both are stated here so that the next measurement, on a beam or on a
corrected injector, starts from them.

Is \bfs worth having? Where a second copy of the data fits, redundancy
already returns correct data and \bfs is at most a complement to it.
Where it does not, on one small medium, \bfs returned correct data or,
once, refused, for 7\,\% of capacity, where ext4 and ext3 returned wrong
data and btrfs could only refuse (\cref{sec:brings}). That is a narrow claim,
bounded by the load measured here, and v4 is planned to widen it or to
find where it stops (\cref{sec:v4}).

\section*{Disclosure}
The analysis scripts and the draft of this report were written with the
assistance of Claude Opus 5.5 (model \texttt{claude-opus-5-5}, Anthropic). Every
measurement was taken by the author's tooling on the author's machines,
and every number in the text is generated from the run records as
described in \apprepro.
